% Extensión estimada: unas 2 páginas, cuatro párrafos y una tabla (ficha 5.5).
% Fuentes: Tablas tab:marco_generadores y tab:marco_sota_reportes y Sección 2.7.3 (sec:estado_arte_reportes) del
% Capítulo 2; Zeng et al., MM-Skin, Tabla 3 (VQA, columna de dermatoscopía, escala 0--100; verificada en el PDF el
% 24/10/2026): LLaVA-Med-7B BLEU-4 8.88, BERTScore 66.60; SkinVL-MM 21.61 y 71.34. Cifras propias de las
% secciones 5.1 a 5.4.
\section{Comparación con el estado del arte}
\label{sec:res_gen_estado_arte}

Como en la clasificación (Sección~\ref{sec:res_estado_arte}), los trabajos de generación de reportes del
Capítulo~\ref{cp:marco} usan otros datos, otras instrucciones y, sobre todo, otras formas de evaluar el texto, de modo
que la comparación ubica el orden de magnitud de los resultados y no ordena los sistemas. La
Tabla~\ref{tab:res_gen_estado_arte} reúne las cifras que tienen un equivalente en este trabajo, en tres bloques: el
parecido con un texto de referencia, la corrección del diagnóstico en el texto y el efecto del contexto recuperado.

\begin{table}[htbp]
    \centering
    \caption{Resultados de la literatura (Capítulo~\ref{cp:marco}) frente a los de este trabajo, en escala de 0 a 1.
    «---»: sin cifra equivalente.}
    \label{tab:res_gen_estado_arte}
    \footnotesize
    \setlength{\tabcolsep}{3pt}
    \begin{tabularx}{\textwidth}{@{}>{\raggedright\arraybackslash}p{2.4cm}
        >{\raggedright\arraybackslash\hspace{0pt}}X
        >{\raggedright\arraybackslash\hspace{0pt}}p{2.5cm}
        >{\centering\arraybackslash}p{1.6cm}>{\centering\arraybackslash}p{2.7cm}@{}}
        \toprule
        \textbf{Trabajo} & \textbf{Datos y protocolo} & \textbf{Métrica} & \textbf{Trabajo citado}
            & \textbf{Este trabajo} \\
        \midrule
        \multicolumn{5}{@{}l}{\textit{Parecido con un texto de referencia}} \\
        \textcite{zeng2025mmskin} & LLaVA-Med sin ajuste; conjunto de prueba de \emph{MM-Skin} de los autores, dermatoscopía,
            preguntas y respuestas & BLEU-4 & 0.089 & 0.001 (LLaVA-Med) \\
        \textcite{zeng2025mmskin} & LLaVA-Med sin ajuste; mismo conjunto de prueba & BERTScore & 0.666 & 0.674 (LLaVA-Med)\newline 0.713 (Mini-VLM) \\
        \textcite{zeng2025mmskin} & SkinVL ajustado con \emph{MM-Skin}; mismo conjunto de prueba & BLEU-4 / BERTScore & 0.216 / 0.713
            & --- \\
        \filasep
        \multicolumn{5}{@{}l}{\textit{Corrección del diagnóstico en el texto}} \\
        \textcite{zhou2024pretrained} & SkinGPT-4; 150 casos reales; cinco dermatólogos
            & Diagnósticos correctos o pertinentes & 0.806 & 0.849 (consistencia, Mini-VLM con Swin-Tiny) \\
        \filasep
        \multicolumn{5}{@{}l}{\textit{Contexto recuperado}} \\
        \textcite{xia2024rule} & LLaVA-Med-1.5 en \emph{MIMIC-CXR}; reportes recuperados agregados sin ajuste
            & Cambio de exactitud al agregar el contexto & $-0.084$ & $+0.013$ (consistencia, Mini-VLM con Swin-Tiny) \\
        \bottomrule
    \end{tabularx}
\end{table}

En el parecido con un texto de referencia, los generadores de este trabajo quedan muy por debajo de los de
\textcite{zeng2025mmskin} en BLEU-4, pero no en BERTScore. El BLEU-4 de LLaVA-Med ajustado con \emph{SkinCaRe} es
casi nulo, mientras que en ese trabajo el mismo modelo sin ajuste comparte con las descripciones de \emph{MM-Skin} algunas frases y
SkinVL, ajustado con ellas, muchas más. En cambio, el BERTScore de LLaVA-Med queda cerca del que obtiene sin ajuste en
ese trabajo, y el de Mini-VLM alcanza el de SkinVL. La comparación no es directa, porque el conjunto de prueba y la instrucción
son otras, pero sugiere que la distancia está más en las palabras que en el contenido: el BLEU-4 premia escribir como el
libro de texto, y solo el modelo ajustado con ese estilo lo hace. Es la misma razón por la que la
Sección~\ref{sec:met_eval_generacion} no usa el Nivel~1 para elegir modelos.

En la corrección del diagnóstico, la cifra más cercana es la de SkinGPT-4~\cite{zhou2024pretrained}, cuyos
dermatólogos consideran correcto o pertinente el 80.6\,\% de los diagnósticos, sobre fotografías de casos reales. La
consistencia de la cascada con Mini-VLM queda en el mismo orden de magnitud, pero mide otra cosa: una búsqueda automática de la
clase verdadera en imágenes dermatoscópicas de siete clases, sin un especialista que juzgue el texto. Del resto de los
trabajos que miden el diagnóstico, DermatoLlama~\cite{yilmaz2025resourceefficient} lo hace sobre datos sintéticos y
\textcite{chakrabarti2025skin}, sobre un único modelo destilado, y ninguno separa el aporte del clasificador del
generador. En este trabajo, en cambio, la cascada muestra que la fidelidad diagnóstica del reporte coincide con la exactitud del clasificador que lo ancla (Sección~\ref{sec:res_gen_fuentes}).

En la recuperación, los resultados coinciden en parte con la advertencia de RULE y MMed-RAG~\cite{xia2024rule,
xia2025mmedrag}. En ellos, el contexto recuperado agregado sin más puede empeorar la respuesta, y \textcite{xia2024rule}
lo miden en \emph{MIMIC-CXR}. Aquí, la recuperación sola también queda por debajo de escribir el voto de los vecinos, de
modo que el generador aprovecha mal el contexto. Sin embargo, junto al diagnóstico del clasificador el contexto no
degrada el reporte, ni siquiera cuando el vecino se equivoca, porque el generador sigue al clasificador
(Sección~\ref{sec:res_gen_rag}). La diferencia es que aquí se recuperan imágenes con su diagnóstico y no reportes de
texto.

Por último, ningún trabajo revisado combina la concordancia entre dos predictores con la generación del reporte. El
ensamble de \textcite{khan2025ensemble} deriva a un especialista los casos en que sus miembros no coinciden, pero no
informa cuántos casos deriva ni con qué acierto, y su texto lo escribe un modelo que no ve la imagen. En este trabajo,
la concordancia separa los casos en que el clasificador acierta nueve de cada diez de aquellos en que acierta la mitad
(Sección~\ref{sec:res_gen_concordancia}), aunque no detecta los errores compartidos. En conjunto, los reportes quedan
en el mismo orden de magnitud que los de la literatura en las métricas comparadas, y la diferencia principal del trabajo está en que
mide por separado el aporte de cada componente del pipeline.
